\documentclass[
  oneside,
  degree=bachelor,
  fullwidthstop=false,
  fontset=none,
  times=false,
  minted=false,
  biblatex=false,
]{tongjithesis}

% ============================================================
%  课程报告通用格式模板
%  说明：
%  1. 这份文件是单文件模板，直接复制后改封面和正文即可。
%  2. 编译方式建议使用 XeLaTeX。
%  3. 这套字体配置已经针对当前这台 mac 做过兼容修复。
% ============================================================

\usepackage{xeCJK}
\setCJKmainfont[
  Path=/System/Library/Fonts/Supplemental/,
  Extension=.ttc,
  FontIndex=3,
  BoldFont=Songti.ttc,
  BoldFeatures={FontIndex=1},
  AutoFakeBold=false
]{Songti.ttc}
\setCJKsansfont[
  Path=/System/Library/Fonts/,
  Extension=.ttc,
  FontIndex=1,
  BoldFont=STHeiti Medium.ttc,
  BoldFeatures={FontIndex=1}
]{STHeiti Light.ttc}
\setCJKmonofont{仿宋_GB2312}[Path=/Users/kinakoxu/Library/Fonts/,Extension=.ttf]
\setCJKfamilyfont{zhsong}[
  Path=/System/Library/Fonts/Supplemental/,
  Extension=.ttc,
  FontIndex=3,
  BoldFont=Songti.ttc,
  BoldFeatures={FontIndex=1},
  AutoFakeBold=false
]{Songti.ttc}
\setCJKfamilyfont{zhhei}[
  Path=/System/Library/Fonts/,
  Extension=.ttc,
  FontIndex=1,
  BoldFont=STHeiti Medium.ttc,
  BoldFeatures={FontIndex=1}
]{STHeiti Light.ttc}
\setCJKfamilyfont{zhkai}{楷体_GB2312}[Path=/Users/kinakoxu/Library/Fonts/,Extension=.ttf]
\setCJKfamilyfont{zhfs}{仿宋_GB2312}[Path=/Users/kinakoxu/Library/Fonts/,Extension=.ttf]
\providecommand{\songti}{\CJKfamily{zhsong}}
\providecommand{\heiti}{\CJKfamily{zhhei}}
\providecommand{\kaishu}{\CJKfamily{zhkai}}
\providecommand{\fangsong}{\CJKfamily{zhfs}}

\begin{document}

% ============================================================
%  封面信息
%  直接修改下面这些字段
% ============================================================
\school{国豪书院}
\major{}
\student{2452627}{徐佳涛}
\thesistitle{课程报告标题第一行\\课程报告标题第二行}{}
\thesistitleeng{Course Report Title in English}{}
\thesisadvisor{张红云}
\thesisdate{\the\year}{\the\month}{\the\day}

\MakeCover
\cleardoublepage

\tableofcontents
\cleardoublepage

\mainmatter

\chapter{问题背景与动态场景建模}
这一章一般写题目来源、研究背景、任务目标、场景抽象方式，以及地图或环境建模方法。

如果是机器人、搜索、机器学习类题目，可以在这里先定义状态空间、动作空间、约束条件和评价指标。

\chapter{项目结构与程序设计}
这一章一般介绍代码结构、主要文件、模块划分、程序执行流程。

可以配合简短代码片段说明主程序如何组织各个算法模块、可视化模块和数据统计模块。

\chapter{算法设计}
这一章写算法原理最合适。可以分小节介绍 BFS、A*、Dijkstra、强化学习、启发式算法等。

建议写法：
\begin{itemize}
  \item 先说明算法核心思想；
  \item 再说明为什么适合本题；
  \item 最后补一两句优缺点。
\end{itemize}

\chapter{核心代码实现}
这一章适合放关键代码，而不是把整份程序贴进来。建议只保留最核心的函数或逻辑。

\begin{lstlisting}[language=Python,caption={示例：A* 搜索框架}]
def a_star_search(grid, start, end):
    open_set = []
    heapq.heappush(open_set, (0, start))
    g_score = {start: 0}

    while open_set:
        _, current = heapq.heappop(open_set)
        if current == end:
            return True
\end{lstlisting}

\chapter{实验结果与动态可视化}
这一章一般写实验环境、测试场景、结果图、动态图、指标表格。

如果需要插图，推荐用下面这种写法：

\begin{figure}[H]
  \centering
  % \includegraphics[width=0.8\textwidth]{your_figure.png}
  \caption{在此处替换为你的实验图标题}
\end{figure}

\chapter{结果分析与讨论}
这一章重点写不同方法之间的比较，包括：
\begin{itemize}
  \item 搜索效率是否更高；
  \item 路径是否更短、更平滑；
  \item 是否容易陷入局部最优；
  \item 是否需要训练；
  \item 算法在复杂场景中的鲁棒性如何。
\end{itemize}

如果题目允许，也可以在这一章写局限性和后续改进方向。

\chapter{总结}
这一章通常用一段到两段话概括全文即可，回答两个问题：
\begin{itemize}
  \item 这次课程设计完成了什么；
  \item 最终得到的主要结论是什么。
\end{itemize}

\end{document}
